% !TEX root = ../diplomarbeit.tex

\section{\emph{Beispiel:} Hardware\authorA}

% Hier beschreibt Autor A seinen Teil der Konzepte.
% Verschiedene Lösungsansätze vorstellen und bewerten.

\subsection{Lösungsansatz 1}

\emph{Hier wird der erste Lösungsansatz detailliert beschrieben, einschließlich der zugrunde liegenden Idee, der theoretischen Grundlagen und der erwarteten Vor- und Nachteile. Grafische Darstellungen wie Diagramme oder Skizzen unterstützen das Verständnis der Konzeptarchitektur.}

\placeholderbox{Beispiel}{%
\textbf{Ansatz: Zentrale Server-basierte Verarbeitung}\\[0.5ex]
Alles wird auf einem Server verarbeitet. Der Client sendet Anfragen an einen zentralen Server, der alle Verarbeitungsschritte durchführt und Ergebnisse zurücksendet. Vorteil: Einfache Wartung und zentrale Kontrolle. Nachteil: Hohe Serverauslastung bei vielen gleichzeitigen Anfragen, Single Point of Failure.%
    %
    \newline\newline
    ...
}

\subsection{Lösungsansatz 2}

\emph{Dieser Abschnitt präsentiert einen alternativen Lösungsansatz mit dessen spezifischen Eigenschaften und Unterschieden zum ersten Konzept. Die Bewertung erfolgt anhand definierter Kriterien wie Machbarkeit, Kosten, Zeitaufwand und technischer Komplexität.}

\placeholderbox{Beispiel}{%
\textbf{Ansatz: Verteilte Client-/Server-Architektur}\\[0.5ex]
Große Datenmengen werden vorab auf den Client übertragen und dort verarbeitet. Der Server bleibt stateless. Vorteil: Bessere Skalierbarkeit, offline-fähig, reduzierte Serverauslastung. Nachteil: Höherer Datentransfer initial, komplexere Client-side Logik, Browser-Speicherlimitierungen.%
    %
    \newline\newline
    ...
}


\section{\emph{Beispiel:} Frontend\authorA}

\subsection{Lösungsansatz 1}

\ldots

\subsection{Lösungsansatz 2}

\ldots